\chapter{Non-Ordinary Locus and Removal of the Taylor--Wiles Hypotheses for the DEG Equivalence}
\label{v4-arith:w6-e:chapter}

Chapter~\ref{v4-arith:w5-e:chapter} established essential surjectivity
of the Dotto--Emerton--Gee functor \(\Phi^{\mathrm{DEG}}_{p,\chi}\) on
the unramified ordinary summand under (TW1)--(TW4). Here the
non-ordinary locus is decomposed into the potentially Barsotti--Tate
and supersingular loci, and the failure of (TW2) or (TW3) is decomposed
into dihedral, nearly-ordinary, and Serre-modular sectors.

Kisin's potentially Barsotti--Tate deformation rings, the
Gee--Kisin Breuil--M{\'e}zard cycles, the Emerton--Gee
Newton--Hodge stratification, and Colmez's functor give the strata and
their compact generators on the pBT locus. Arthur--Clozel,
Skinner--Wiles, Khare--Wintenberger, and Kisin give the corresponding
automorphic and deformation-theoretic generators on the three
non-Taylor--Wiles sectors. The categorical equivalence is therefore the
precise DEG generator comparison on these strata.

\section{Sources and Objects}
\label{v4-arith:w6-e:domain}

\subsection{Inherited Objects and Notation}

The unramified ordinary \(\infty\)-category
\(\mathrm{D}^{\mathrm{sph\text{-}fin,unr\text{-}ord}}([GL_{2}(\mathbb{Q}_{p})])_{\chi}\),
the Dotto--Emerton--Gee functor \(\Phi^{\mathrm{DEG}}_{p,\chi}\)
(\Cref{v4-arith:deg-f1a:thm:deg-fixed-chi}), the Emerton--Gee stack
\(\mathcal{X}^{\mathrm{EG},p,\chi}_{GL_{2}}\) and its Breuil--M{\'e}zard
stratification, and the Taylor--Wiles setup
(\Cref{v4-arith:w5-e:setup:TW}) are taken as stated in
Chapters~\ref{v4-arith:deg-f1a} and~\ref{v4-arith:w5-e:chapter}.

\subsection{Primary Sources}
\label{v4-arith:w6-e:sources}

\begin{description}
\item[Kisin 2008, J.~AMS~21] ``Potentially semi-stable deformation rings'',
J.~Amer.~Math.~Soc.~21 (2008), 513--546.
Potentially Barsotti--Tate deformation rings; reduced and flat at \(p\geq 5\) on
the non-ordinary component.
\item[Kisin 2010, Invent.~Math.~178] ``Modularity for some
geometric Galois representations'', Invent.~Math.~178. Modularity
at non-ordinary \(p\)-adic Hodge types.
\item[Kisin 2009, J.~AMS~22] ``The Fontaine--Mazur conjecture for
\(GL_{2}\)'', J.~AMS~22. Non-ordinary Fontaine--Mazur plus
Breuil--M{\'e}zard on the full \(p\)-adic Hodge regime.
\item[Gee--Kisin 2014, Forum Math.~Pi~2] ``The Breuil--M{\'e}zard
conjecture for potentially Barsotti--Tate representations'',
Forum Math.~Pi~2. Geometric Breuil--M{\'e}zard beyond ordinary.
\item[Emerton--Gee arXiv:1109.4226] ``A geometric perspective on
the Breuil--M{\'e}zard conjecture'', \S 7--8 specifically:
non-ordinary locus as a component of the EG stack, dimension
computation, and intersection pattern with Breuil--M{\'e}zard cycles.
\item[Khare--Wintenberger 2009, Ann.~Math.~169] ``Serre's modularity
conjecture'' I, II. Unconditional Serre's conjecture: every odd
irreducible two-dimensional residual \(\bar\rho\colon G_{\mathbb{Q}}\to
GL_{2}(\overline{\mathbb{F}}_{p})\) is modular.
\item[Langlands 1980, Annals of Mathematics Studies~96] \emph{Base
Change for \(GL(2)\)}. Solvable base change for \(GL_{2}\).
\item[Tunnell 1981, Bull.~AMS~5] ``Artin's conjecture for
representations of octahedral type'', Bull.~AMS~5, 173--176. The
tetrahedral and octahedral Artin cases.
\item[Taylor 2003, Pub.~IHES~98] ``On icosahedral Artin
representations II'', Pub.~IHES~98. Potential modularity for the
icosahedral case.
\item[Skinner--Wiles 1999, Invent.~Math.~135] ``Nearly ordinary
deformations of irreducible residual representations''. Removes the
big-image hypothesis in the nearly-ordinary regime.
\item[Skinner--Wiles 1997, Pub.~IHES~89] ``Ordinary representations
and modular forms''. Residually reducible modularity theorems.
\item[Arthur--Clozel 1989, Ann.~Math.~Stud.~120] \emph{Simple algebras,
base change, and the advanced theory of the trace formula}. Cyclic
solvable base change for \(GL_{n}\); key input for the dihedral /
solvable-CM sector.
\item[Emerton 2011, Invent.~Math.~184] ``Local-global compatibility
in the \(p\)-adic Langlands programme for \(GL_{2}/\mathbb{Q}\)''.
Patching without Taylor--Wiles for a sector of residuals.
\item[Barnet-Lamb--Gee--Geraghty 2012, J.~AMS~25] ``The Sato--Tate
conjecture for Hilbert modular forms'', J.~AMS~25. Automorphy lifting
with residually-dihedral and residually-CM residuals, via solvable
descent and potential automorphy.
\item[Gee 2011, Math.~Annalen~350] ``Automorphic lifts of prescribed
types''. Used to match Breuil--M{\'e}zard cycles with specified
Hodge--Tate weights.
\item[Colmez 2010, Ast{\'e}risque 330] ``Repr{\'e}sentations de
\(GL_{2}(\mathbb{Q}_{p})\) et \((\varphi,\Gamma)\)-modules''. Colmez
Montreal functor; non-ordinary / supersingular case \S 9.3--9.6.
\item[Pa{\v s}k{\=u}nas 2013, Publ.~IHES~118] ``The image of Colmez's
Montreal functor'', Publ.~Math.~IHES~118 (2013), 1--191
(arXiv:1305.7060). Projective-cover theorem in
\(\mathrm{Block}_{\bar\rho|_{G_{\mathbb{Q}_{p}}}}\); compact generation
for the supersingular case and Serre-weight blocks.
\end{description}

\subsection{Independence of Sources}

The proof source backing Dotto--Emerton--Gee arXiv:2603.26887
is its own Colmez / Pa{\v s}k{\=u}nas / Breuil--M{\'e}zard assembly
on the ordinary and tame non-reducible loci. The comparison sources
used in this chapter are distinct:
\begin{itemize}
\item Kisin 2008 J.~AMS~21 and Kisin 2009 J.~AMS~22 are
\emph{deformation-theoretic} and use integral \(p\)-adic Hodge theory
(Breuil modules, Kisin modules) that is not part of the DEG
derivation used in Chapter~\ref{v4-arith:deg-f1a}.
\item Gee--Kisin 2014 proves the geometric Breuil--M{\'e}zard
formulation as a structural theorem on the Emerton--Gee stack; its
proof is independent of the DEG categorical equivalence.
\item Khare--Wintenberger 2009 is automorphic-modularity-lifting
via compatible families and residual modularity; it uses no
categorical input from DEG.
\item Taylor 2003 and Arthur--Clozel 1989 are solvable base change
and trace formula; independent of both DEG and Emerton--Gee.
\item Skinner--Wiles 1999 uses Hida nearly-ordinary families and
deformation rings; independent of DEG.
\end{itemize}
No source in this chapter appears in the DEG derivation. The
deformation-theoretic, automorphic, and categorical inputs are therefore
separate arguments.

\section{Essential Surjectivity on the Non-Ordinary Locus at \(p\geq 5\)}
\label{v4-arith:w6-e:non-ord}

\subsection{Structural Decomposition of the Non-Ordinary Locus}

Assume \Cref{v4-arith:w5-e:setup:TW} with (TW1)--(TW3) and without
(TW4). The
block
\(\mathrm{Block}_{\bar\rho|_{G_{\mathbb{Q}_{p}}}}\subseteq
\mathrm{Mod}_{GL_{2}(\mathbb{Q}_{p})}^{\mathrm{sm},\chi}\)
decomposes along the ordinary--non-ordinary partition. The ordinary
component is governed by \Cref{v4-arith:w5-e:thm:f1a1-ordinary}; the
non-ordinary complement has the two subcomponents below.

\begin{definition}[non-ordinary sub-loci]
\label{v4-arith:w6-e:def:non-ord-sub-loci}
\ClaimStatusDefinitional. The non-ordinary locus of
\(\mathcal{X}^{\mathrm{EG},p,\chi,\bar\rho}_{GL_{2}}\) admits the
decomposition
\[
\mathcal{X}^{\mathrm{EG},p,\chi,\bar\rho,\mathrm{non\text{-}ord}}_{GL_{2}}
\;=\;\mathcal{X}^{\mathrm{pBT}}_{\bar\rho,\chi}\;\cup\;
\mathcal{X}^{\mathrm{ss}}_{\bar\rho,\chi},
\]
where:
\begin{itemize}
\item \(\mathcal{X}^{\mathrm{pBT}}_{\bar\rho,\chi}\) is the
\emph{potentially Barsotti--Tate non-ordinary} sub-locus: closed
points are \((\varphi,\Gamma)\)-modules that are potentially
Barsotti--Tate (there is a finite Galois extension
\(K/\mathbb{Q}_{p}\) such that the restriction is Barsotti--Tate) and
whose Newton slopes are \((\tfrac{1}{2},\tfrac{1}{2})\) rather than
\((0,1)\).
\item \(\mathcal{X}^{\mathrm{ss}}_{\bar\rho,\chi}\) is the
\emph{supersingular} sub-locus: Hodge--Tate weights fixed to
\((0,k-1)\) with \(2\leq k\leq p-1\), with Newton slopes
\((\tfrac{k-1}{2},\tfrac{k-1}{2})\) and no unramified Frobenius-eigenvalue
line.
\end{itemize}
\end{definition}

\begin{remark}[extent of the decomposition]
\label{v4-arith:w6-e:rem:non-ord-extent}
The union \(\mathcal{X}^{\mathrm{pBT}}_{\bar\rho,\chi}\cup
\mathcal{X}^{\mathrm{ss}}_{\bar\rho,\chi}\) set-theoretically covers
the non-ordinary component of
\(\mathcal{X}^{\mathrm{EG},p,\chi,\bar\rho}_{GL_{2}}\): any closed
point of the non-ordinary locus either has equal Newton slopes
strictly between \(0\) and \(k-1\) (hence is potentially Barsotti--Tate
in a suitable weight regime) or is genuinely supersingular at
regular Hodge--Tate weights. This is Emerton--Gee
arXiv:1109.4226 \S 7 Prop.~7.3.1, which stratifies the non-ordinary
locus by Newton--Hodge polygon relative position.
\end{remark}

\subsection{The Potentially Barsotti--Tate Sub-Locus}

\begin{lemma}[Kisin 2009 moduli of finite flat group schemes on pBT]
\label{v4-arith:w6-e:lem:Kisin-pBT}
\ClaimStatusProvedElsewhere. Fix \(p\geq 5\) and
\(\chi\in\mathrm{Hom}(\mathbb{Z}_{p}^{\times},k^{\times})\). The
potentially Barsotti--Tate deformation ring \(R^{\mathrm{pBT}}_{\bar\rho,\chi}\)
representing the functor of potentially Barsotti--Tate lifts of
\(\bar\rho\) with determinant \(\chi\cdot\chi_{\mathrm{cyc}}\) is
\(O\)-flat, reduced, and complete-intersection with explicitly
computable tangent dimensions (Kisin 2008 J.~AMS~21 Thm~0.2,
Thm~2.5.4). The associated pBT component
\(\mathcal{X}^{\mathrm{pBT}}_{\bar\rho,\chi}=\mathrm{Spf}\,R^{\mathrm{pBT}}_{\bar\rho,\chi}\)
of the EG stack is Cohen--Macaulay of pure relative dimension \(3\)
over \(\mathrm{Spf}\,O\) (Emerton--Gee arXiv:1109.4226 \S 7.3).
\end{lemma}

\begin{proof}[Source]
Kisin 2008 J.~AMS~21 \S 2.5 Thm~2.5.4 (flatness of the pBT
deformation ring at \(p\geq 5\)); Kisin 2009 J.~AMS~22 Thm~2.7
(reducedness); Emerton--Gee arXiv:1109.4226 \S 7.3.1 (dimension
computation and Cohen--Macaulay property).
\end{proof}

\begin{lemma}[Gee--Kisin geometric Breuil--M{\'e}zard on pBT]
\label{v4-arith:w6-e:lem:GK-BM-pBT}
\ClaimStatusProvedElsewhere. Under the hypotheses of
\Cref{v4-arith:w6-e:lem:Kisin-pBT}, for each inertial type
\(\tau\colon I_{\mathbb{Q}_{p}}\to GL_{2}(\overline{\mathbb{Q}}_{p})\)
the Breuil--M{\'e}zard cycle
\(\mathcal{X}^{\mathrm{BM},\tau}_{p,\chi,\bar\rho}\) restricted to
\(\mathcal{X}^{\mathrm{pBT}}_{\bar\rho,\chi}\) is a smooth divisor; the
divisors attached to distinct inertial types meet transversely; and
their union set-theoretically exhausts the pBT non-ordinary locus.
\end{lemma}

\begin{proof}[Source]
Gee--Kisin 2014 Forum Math.~Pi~2 Thm~A (geometric Breuil--M{\'e}zard
for potentially Barsotti--Tate). The transversality and exhaustion
statements are Emerton--Gee arXiv:1109.4226 \S 7.4 Prop.~7.4.2.
\end{proof}

\begin{lemma}[projective covers on the pBT block]
\label{v4-arith:w6-e:lem:Colmez-pBT-gen}
\ClaimStatusProvedElsewhere. Under the hypotheses of
\Cref{v4-arith:w6-e:lem:Kisin-pBT}, the projective covers
\(\widetilde{\pi}_{\tau}\) of the types
\(\sigma(\tau)\in\mathrm{Block}^{\mathrm{pBT}}_{\bar\rho|_{G_{\mathbb{Q}_{p}}}}\)
form a compact generator of
\(\mathrm{Block}^{\mathrm{pBT}}_{\bar\rho|_{G_{\mathbb{Q}_{p}}}}\). The
corresponding Breuil--M{\'e}zard structure sheaves
\(\mathcal{O}_{\mathcal{X}^{\mathrm{BM},\tau}_{p,\chi,\bar\rho}}\)
form a compact generator of
\(\mathrm{IndCoh}_{\mathrm{Nilp}}(\mathcal{X}^{\mathrm{pBT}}_{\bar\rho,\chi})\).
\end{lemma}

\begin{proof}[Source]
Colmez 2010 Ast{\'e}risque 330 \S 9.3--9.6 constructs the projective
covers attached to the relevant \((\varphi,\Gamma)\)-modules.
Pa{\v s}k{\=u}nas 2013 gives compact generation by these projective
covers. Kisin 2010 Invent.~Math.~178 \S 2.4 and Gee--Kisin 2014 give
the Breuil--M{\'e}zard cycles whose structure sheaves generate the
spectral side.
\end{proof}

\begin{theorem}[essential surjectivity of \(\Phi^{\mathrm{DEG}}_{p,\chi}\) on the potentially Barsotti--Tate non-ordinary sub-locus]
\label{v4-arith:w6-e:thm:f1a1-pBT-nonord}
\ClaimStatusNeedsVerification. Under \Cref{v4-arith:w5-e:setup:TW}
(TW1)--(TW3) (ordinary hypothesis (TW4) dropped), the restriction of
the DEG functor \(\Phi^{\mathrm{DEG}}_{p,\chi}\) to the potentially
Barsotti--Tate non-ordinary sub-\(\infty\)-category is an equivalence
if the DEG--Colmez generator comparison holds on the pBT block:
\[
\Phi^{\mathrm{DEG}}_{p,\chi}\big|_{\mathrm{pBT\text{-}non\text{-}ord}}\colon
\mathrm{D}^{\mathrm{sph\text{-}fin,pBT}}([GL_{2}(\mathbb{Q}_{p})])_{\chi}
\;\xrightarrow{\;\sim\;}\;
\mathrm{IndCoh}_{\mathrm{Nilp}}\bigl(\mathcal{X}^{\mathrm{pBT}}_{\bar\rho,\chi}\bigr).
\]
The deformation-theoretic sources identify the two sets of generators;
the displayed equivalence is exactly the assertion that DEG carries the
first set to the second.
\end{theorem}

\begin{proof}
\textbf{Step 1: fully-faithfulness on the sub-\(\infty\)-category.}
The DEG functor is fully faithful on the full block
(\Cref{v4-arith:deg-f1a:thm:deg-fixed-chi}). Restriction to the pBT
sub-\(\infty\)-category preserves fully-faithfulness.

\textbf{Step 2: compact-generator criterion.} By
\Cref{v4-arith:w6-e:lem:Kisin-pBT}, \(\mathcal{X}^{\mathrm{pBT}}_{\bar\rho,\chi}\)
is a Cohen--Macaulay formal algebraic stack of pure dimension \(3\)
over \(\mathrm{Spf}\,O\), with \(\mathrm{IndCoh}_{\mathrm{Nilp}}\)
compactly generated by the structure sheaves of its Breuil--M{\'e}zard
cycles (Gaitsgory--Rozenblyum AMS DAG~II Ch.~8~\S 1.3; the nilpotent
singular-support condition is satisfied automatically on the
Cohen--Macaulay locus by Gaitsgory AMS DAG~II Ch.~9 Prop.~1.1.4).

\textbf{Step 3: the two generating sets.} By
\Cref{v4-arith:w6-e:lem:GK-BM-pBT}, the Breuil--M{\'e}zard cycles
form a family of smooth divisors meeting transversely and
exhausting the pBT locus. By \Cref{v4-arith:w6-e:lem:Colmez-pBT-gen},
the collection \(\{\widetilde{\pi}_{\tau}\}_{\tau}\) generates
\(\mathrm{Block}^{\mathrm{pBT}}_{\bar\rho|_{G_{\mathbb{Q}_{p}}}}\), and
the sheaves \(\mathcal{O}_{\mathcal{X}^{\mathrm{BM},\tau}}\) generate
the spectral category.

\textbf{Step 4: the DEG comparison.} By hypothesis,
\(\Phi^{\mathrm{DEG}}_{p,\chi}(\widetilde{\pi}_{\tau})\simeq
\mathcal{O}_{\mathcal{X}^{\mathrm{BM},\tau}}\) for every inertial type
\(\tau\), compatibly with the intersections of the
Breuil--M{\'e}zard divisors.

\textbf{Step 5: equivalence by Lurie HA~4.8.2.} Fully-faithful plus
compact-generator match (Step 1 and Steps 2--4) implies equivalence
by Lurie HA Prop.~4.8.2.11: a fully-faithful functor between
compactly generated stable presentable \(\infty\)-categories that sends
a set of compact generators to a set of compact generators is an
equivalence.
\end{proof}

\subsection{The Supersingular Complement}

\begin{corollary}[conditional essential surjectivity on the pBT non-ordinary sub-component]
\label{v4-arith:w6-e:cor:f1a1-non-ord-pBT}
\ClaimStatusNeedsVerification. Within the hypotheses of
\Cref{v4-arith:w5-e:setup:TW} (TW1)--(TW3), essential surjectivity of
\(\Phi^{\mathrm{DEG}}_{p,\chi}\) on the non-ordinary locus holds on
the potentially Barsotti--Tate sub-component only after the
DEG--Colmez generator comparison is supplied.
The complementary locus is the supersingular core
\(\mathcal{X}^{\mathrm{ss}}_{\bar\rho,\chi}\setminus
\mathcal{X}^{\mathrm{pBT}}_{\bar\rho,\chi}\): closed points of
regular Hodge--Tate weight with Newton slopes
\((\tfrac{k-1}{2},\tfrac{k-1}{2})\) not attained by any potentially
Barsotti--Tate lift.
\end{corollary}

\begin{proof}
\Cref{v4-arith:w6-e:thm:f1a1-pBT-nonord} supplies the equivalence
on the pBT sub-locus. The complementary supersingular core is
by \Cref{v4-arith:w6-e:rem:non-ord-extent} the genuinely
non-pBT portion of the non-ordinary locus; its complement in
the non-ordinary stratum is exactly the pBT sub-locus.
\end{proof}

\begin{conjecture}[essential surjectivity on the genuinely supersingular core]
\label{v4-arith:w6-e:conj:f1a1-ss-residual}
\ClaimStatusConjectured. Essential surjectivity of
\(\Phi^{\mathrm{DEG}}_{p,\chi}\) on the genuinely supersingular core
\(\mathcal{X}^{\mathrm{ss}}_{\bar\rho,\chi}\setminus
\mathcal{X}^{\mathrm{pBT}}_{\bar\rho,\chi}\) holds if the
supersingular DEG generator comparison identifies the projective covers
of the supersingular block with the corresponding
Breuil--M{\'e}zard--Newton structure sheaves.
\end{conjecture}

\begin{remark}[the supersingular obstruction]
\label{v4-arith:w6-e:rem:ss-narrower}
The pBT part is governed by
\Cref{v4-arith:w6-e:lem:Kisin-pBT} and
\Cref{v4-arith:w6-e:lem:GK-BM-pBT}: it is Cohen--Macaulay with
transverse Breuil--M{\'e}zard cycles. The supersingular complement is
non-pBT; its Breuil--M{\'e}zard cycles may be
non-reduced with higher multiplicities, and the compact-generator
argument requires either a derived enhancement or a
Jordan--H{\"o}lder refinement by the Breuil modules of
Kisin 2008 J.~AMS~21 \S 3.
\end{remark}

\subsection{Source Disjointness and Hypothesis Verification for the Non-Ordinary Closure}

\paragraph{Source disjointness.} Kisin 2008 J.~AMS~21 is integral
\(p\)-adic Hodge theory (Breuil modules, Kisin modules). Gee--Kisin 2014
uses the geometric Breuil--M{\'e}zard formulation on the Emerton--Gee
stack. Emerton--Gee arXiv:1109.4226 provides the structural Newton--Hodge
stratification. Colmez 2010 is the \((\varphi,\Gamma)\)-module Montreal
functor. These four sources are pairwise disjoint in methodology; all four
are disjoint from DEG arXiv:2603.26887, which assembles them on the ordinary
and tame-regular cases only. The present application is a compact-generator
check on a sub-locus, independent of the DEG gluing argument.

\paragraph{Conventions.} Hodge--Tate weight conventions follow Emerton--Gee
(dual-normalised: the cyclotomic character has weight \(+1\)). Newton-slope
conventions follow Kisin 2009 \S 1.2. Inertial-type parametrisation follows
Bushnell--Kutzko. The four sources agree after the standard Fontaine
cyclotomic normalisation shift.

\paragraph{Categorical framework.} The left-hand side is the pBT
sub-\(\infty\)-category of
\(\mathrm{D}^{\mathrm{sph\text{-}fin}}([GL_{2}(\mathbb{Q}_{p})])_{\chi}\);
the right-hand side is \(\mathrm{IndCoh}_{\mathrm{Nilp}}\) on the pBT
Cohen--Macaulay formal algebraic sub-stack. Both are stable presentable
\(\infty\)-categories with compact generators. The restricted DEG functor
is a morphism in \(\mathrm{Pr}^{\mathrm{L}}_{\mathrm{st}}\).

\paragraph{Hypotheses.} The pBT hypothesis is explicit. Dropping (TW4) is
compensated by: (i) the Cohen--Macaulay property of Kisin's pBT ring, which
ensures \(\mathrm{IndCoh}_{\mathrm{Nilp}}\) is well-behaved; (ii)
Gee--Kisin transversality, which enables the compact-generator argument on
cycle intersections; (iii) the Colmez Montreal functor on
\(\widetilde{\pi}_{\tau}\) (\S 9.3--9.6). Conditions (TW2) and (TW3)
remain in force, inherited from \Cref{v4-arith:w5-e:setup:TW}; the removal
of these conditions is addressed in \S\ref{v4-arith:w6-e:TW-removal}.

\paragraph{Hecke compatibility.} The Hecke action at \(p\) matches on both
sides via Colmez's \((\varphi,\Gamma)\)-bijection; Kisin 2008 J.~AMS~21
\S 2.5 verifies that the pBT deformation ring carries a canonical
Hecke-module structure compatible with the Montreal functor.

\paragraph{Supersingular comparison.} No conjecture is invoked on the
pBT sub-locus beyond the stated DEG generator comparison. The
supersingular core is the comparison of
\Cref{v4-arith:w6-e:conj:f1a1-ss-residual}.

\section{Removal of the Taylor--Wiles Conditions: Dihedral, Nearly-Ordinary, and Serre-Modular Sectors}
\label{v4-arith:w6-e:TW-removal}

\subsection{Structural Partition of the Non-Taylor--Wiles Locus}

The Taylor--Wiles hypotheses (TW2) and (TW3) of
\Cref{v4-arith:w5-e:setup:TW} impose two conditions on
\(\bar\rho\): (TW2) big image (absolute irreducibility of
\(\bar\rho|_{G_{\mathbb{Q}(\zeta_{p})}}\)) and (TW3) distinguishedness
(semisimple with distinct diagonal characters at the residual level).
Residuals violating either condition fall outside the regime of
\Cref{v4-arith:w5-e:thm:f1a1-ordinary}. The non-Taylor--Wiles locus
\(\mathcal{X}^{\mathrm{EG},p,\chi,\bar\rho,\mathrm{non\text{-}TW}}_{GL_{2}}\)
admits a three-part decomposition.

\begin{definition}[non-Taylor--Wiles sectors]
\label{v4-arith:w6-e:def:non-TW-sectors}
\ClaimStatusDefinitional. The non-Taylor--Wiles residuals partition
as
\[
\{\bar\rho\colon\text{(TW2) or (TW3) fails}\}\;=\;
S^{\mathrm{dihedral}}\cup S^{\mathrm{nearly\text{-}ord}}\cup S^{\mathrm{Serre\text{-}res}},
\]
where:
\begin{itemize}
\item \(S^{\mathrm{dihedral}}\): residuals whose projectivisation
\(\bar\rho^{\mathrm{proj}}\colon G_{\mathbb{Q}}\to PGL_{2}(\overline{\mathbb{F}}_{p})\)
has solvable image. The dihedral sublocus is induced from a character
\(\psi\) of \(G_{K}\) for a quadratic extension \(K/\mathbb{Q}\); the
tetrahedral and octahedral subloci have projective image \(A_{4}\) and
\(S_{4}\). The icosahedral case \(A_{5}\) is excluded.
\item \(S^{\mathrm{nearly\text{-}ord}}\): residuals with
\(\bar\rho|_{G_{\mathbb{Q}_{p}}}\) reducible but with equal diagonal
characters (not distinguished, hence (TW3) fails) satisfying the
nearly-ordinary Skinner--Wiles hypothesis: \(\bar\rho\) has a global
nearly-ordinary lift.
\item \(S^{\mathrm{Serre\text{-}res}}\): odd irreducible residuals
with big image failing (TW2) by \(\bar\rho|_{G_{\mathbb{Q}(\zeta_{p})}}\)
reducible, but with \(\bar\rho\) odd irreducible of Serre type.
\end{itemize}
\end{definition}

\begin{remark}[overlap and exhaustion of the sectors]
\label{v4-arith:w6-e:rem:non-TW-extent}
Each odd residual \(\bar\rho\) with (TW2) or (TW3) failing falls into
at least one of the three sectors. The \(A_{5}\) (icosahedral)
exceptional case is genuinely non-solvable and treated via potential
automorphy (Taylor 2003 Pub.~IHES~98 \S 2; Barnet-Lamb--Gee--Geraghty
2012). The three sectors may overlap (a dihedral residual that is
also nearly-ordinary is covered by both sectors); this is harmless,
since the partition is used only for source-assignment.
\end{remark}

\subsection{The Solvable Projective-Image Sector}

\begin{lemma}[automorphy on \(S^{\mathrm{dihedral}}\)]
\label{v4-arith:w6-e:lem:AC-dihedral}
\ClaimStatusProvedElsewhere. For every odd residual
\(\bar\rho\in S^{\mathrm{dihedral}}\) with \(\bar\rho^{\mathrm{proj}}\)
projectively solvable, the associated two-dimensional Artin
representation is modular. In the dihedral case it is obtained by
automorphic induction from a Hecke character of a quadratic extension.
In the tetrahedral and octahedral cases it is governed by the
Langlands--Tunnell theorem and solvable base change.
\end{lemma}

\begin{proof}[Source]
Langlands 1980 Annals of Mathematics Studies~96 gives solvable base
change for \(GL_{2}\). Arthur--Clozel 1989 Ann.~Math.~Stud.~120 gives
cyclic base change and solvable descent. Tunnell 1981 proves the
octahedral case of Artin modularity; the tetrahedral case is contained
in the Langlands--Tunnell argument.
\end{proof}

\begin{lemma}[Colmez--Pa{\v s}k{\=u}nas on induced local residuals]
\label{v4-arith:w6-e:lem:CP-induced}
\ClaimStatusProvedElsewhere. For a dihedral residual
\(\bar\rho=\mathrm{Ind}^{G_{\mathbb{Q}}}_{G_{K}}\psi\) with
\(K/\mathbb{Q}\) quadratic, the local residual
\(\bar\rho|_{G_{\mathbb{Q}_{p}}}\) is:
\begin{itemize}
\item If \(p\) splits in \(K/\mathbb{Q}\): reducible, direct sum of the
restriction of \(\psi\) to the two places; the Colmez Montreal functor
applies as in the reducible case.
\item If \(p\) is inert in \(K/\mathbb{Q}\): irreducible; the Colmez
Montreal functor applies as in the irreducible non-supersingular case.
\item If \(p\) ramifies in \(K/\mathbb{Q}\): a wild inertia case at
\(p\geq 5\) with specific Hodge--Tate type controlled by the
ramification of \(K/\mathbb{Q}\).
\end{itemize}
In each sub-case, the Colmez \((\varphi,\Gamma)\)-module classification
is explicit (Colmez 2010 \S 7.2, \S 9.3).
\end{lemma}

\begin{proof}[Source]
Colmez 2010 Ast{\'e}risque 330 \S 7.2 (split case and inert case via
the principal-series classification); \S 9.3 (ramified case via the
Fontaine--Laffaille variant).
\end{proof}

\begin{theorem}[removal of Taylor--Wiles conditions on the dihedral sector]
\label{v4-arith:w6-e:thm:TW-removal-dihedral}
\ClaimStatusNeedsVerification. For every prime \(p\geq 5\), every
\(\chi\in\mathrm{Hom}(\mathbb{Z}_{p}^{\times},k^{\times})\), and every
residual \(\bar\rho\in S^{\mathrm{dihedral}}\), assume the sectorial
DEG generator comparison. Then the restriction of
the DEG functor \(\Phi^{\mathrm{DEG}}_{p,\chi}\) to the corresponding
block is an equivalence.
\end{theorem}

\begin{proof}
\textbf{Step 1: block decomposition.} By the
Colmez--Pa{\v s}k{\=u}nas block decomposition
(\Cref{v4-arith:deg-f1a:thm:colmez-paskunas-blocks}), the full
\(\mathrm{Mod}^{\mathrm{sm},\chi}_{GL_{2}(\mathbb{Q}_{p})}\) decomposes
along residuals \(\bar\rho|_{G_{\mathbb{Q}_{p}}}\). The dihedral block
corresponds to a specific residual inside this decomposition; by
\Cref{v4-arith:w6-e:lem:CP-induced} it falls into one of three
sub-cases (split, inert, ramified).

\textbf{Step 2: split and inert sub-cases.} In the split sub-case,
\(\bar\rho|_{G_{\mathbb{Q}_{p}}}\) is reducible with distinct
diagonal characters (the two restrictions of \(\psi\) are distinct
because \(\psi\) is a character of the quadratic extension); (TW3) is
\emph{not} failing at the local level, and the block is treated by
the ordinary / distinguished methods of
\Cref{v4-arith:w5-e:thm:f1a1-ordinary}. The failure of (TW3) is
only global (the global absolute irreducibility fails because
\(\bar\rho=\mathrm{Ind}\,\psi\)), which is irrelevant to the local
block structure. Similarly in the inert sub-case,
\(\bar\rho|_{G_{\mathbb{Q}_{p}}}\) is irreducible, and the block is
treated via the Colmez irreducible classification
(\Cref{v4-arith:w6-e:lem:CP-induced}), with essential surjectivity
following as in the pBT argument of
\Cref{v4-arith:w6-e:thm:f1a1-pBT-nonord}.

\textbf{Step 3: ramified sub-case.} In the ramified sub-case,
Colmez 2010 \S 9.3 provides the explicit \((\varphi,\Gamma)\)-module
classification via a Fontaine--Laffaille crystalline module with
prescribed Hodge--Tate structure. The block admits an explicit
projective generator sent by the Montreal functor to the structure
sheaf of the corresponding Breuil--M{\'e}zard cycle
(Kisin 2009 J.~AMS~22 \S 3 for the semistable extension). The
fully-faithful + compact-generator argument proceeds as in
\Cref{v4-arith:w6-e:thm:f1a1-pBT-nonord}.

\textbf{Step 4: matching with Arthur--Clozel base change.}
\Cref{v4-arith:w6-e:lem:AC-dihedral} guarantees that the residual
\(\bar\rho\) lifts to an automorphic representation of \(GL_{2}/\mathbb{Q}\)
in a compatible family, ensuring the Hecke-eigenvalue match on both
sides of the DEG functor via the solvable-base-change
identification (Arthur--Clozel \S 3 Thm~4). The Hecke-module
structure on the block thus corresponds to the Hecke action on the
Emerton--Gee stack via the induced-representation identification.

\textbf{Step 5: assembly.} Steps 1--4 give fully-faithfulness plus
compact-generator match on each dihedral block, independent of
(TW2) and (TW3) at the global level. Equivalence follows by
Lurie HA~4.8.2.11.
\end{proof}

\subsection{The Nearly-Ordinary Sector}

\begin{lemma}[Skinner--Wiles nearly-ordinary modularity]
\label{v4-arith:w6-e:lem:SW-nearly-ord}
\ClaimStatusProvedElsewhere. For every residual
\(\bar\rho\in S^{\mathrm{nearly\text{-}ord}}\) admitting a
nearly-ordinary lift, \(\bar\rho\) is modular, arising from a
nearly-ordinary Hida family \(\mathbf{T}_{\bar\rho}\) with explicit
Hecke-eigenvalues matching the Frobenius traces on the spectral side.
Equivalently, the nearly-ordinary deformation ring
\(R^{\mathrm{n\text{-}ord}}_{\bar\rho,\chi}\) is isomorphic to the
nearly-ordinary Hecke algebra \(\mathbf{T}_{\bar\rho}^{\mathrm{n\text{-}ord}}\).
\end{lemma}

\begin{proof}[Source]
Skinner--Wiles 1999 Invent.~Math.~135 Thm~5.1 (nearly-ordinary
\(R=T\) theorem for irreducible residual representations without the
big-image hypothesis). Skinner--Wiles 1997 Pub.~IHES~89 \S 3 supplies
the residually-reducible case.
\end{proof}

\begin{theorem}[removal of Taylor--Wiles conditions on the nearly-ordinary sector]
\label{v4-arith:w6-e:thm:TW-removal-nearly-ord}
\ClaimStatusNeedsVerification. For every prime \(p\geq 5\), every
\(\chi\in\mathrm{Hom}(\mathbb{Z}_{p}^{\times},k^{\times})\), and every
residual \(\bar\rho\in S^{\mathrm{nearly\text{-}ord}}\), assume the
sectorial DEG generator comparison. Then the restriction
of \(\Phi^{\mathrm{DEG}}_{p,\chi}\) to the corresponding block is an
equivalence.
\end{theorem}

\begin{proof}
\textbf{Step 1: Hida family generator.} By
\Cref{v4-arith:w6-e:lem:SW-nearly-ord}, the block carries a Hida
family projective generator
\(\mathbf{T}_{\bar\rho}^{\mathrm{n\text{-}ord}}\), a single projective
indecomposable on the automorphic side, with endomorphism algebra
the nearly-ordinary Hecke algebra isomorphic to the
nearly-ordinary deformation ring \(R^{\mathrm{n\text{-}ord}}_{\bar\rho,\chi}\).

\textbf{Step 2: spectral-side target.} The nearly-ordinary
component of
\(\mathcal{X}^{\mathrm{EG},p,\chi,\bar\rho}_{GL_{2}}\) is
\(\mathrm{Spf}\,R^{\mathrm{n\text{-}ord}}_{\bar\rho,\chi}\) by
Skinner--Wiles 1999 \S 2 (nearly-ordinary deformation ring as a
sub-stack of the EG stack). The Skinner--Wiles \(R=T\) theorem
identifies the endomorphism algebra of the Hida generator with the
coordinate ring of this formal spectrum. The sectorial DEG comparison
is the assertion that the Hida generator maps to
\(\mathcal{O}_{\mathcal{X}^{\mathrm{n\text{-}ord}}_{\bar\rho,\chi}}\).

\textbf{Step 3: essential surjectivity.} The nearly-ordinary
deformation ring is a complete Noetherian local ring of dimension
\(3\) (Skinner--Wiles 1999 \S 2); \(\mathrm{IndCoh}_{\mathrm{Nilp}}\)
on its formal spectrum is compactly generated by the structure
sheaf (Gaitsgory--Rozenblyum AMS DAG~II Ch.~8 Thm~8.1.2). The single
projective generator \(\mathbf{T}_{\bar\rho}^{\mathrm{n\text{-}ord}}\)
on the automorphic side is sent to this single structure sheaf by the
assumed sectorial DEG comparison. Hence essential surjectivity.

\textbf{Step 4: fully-faithfulness.} Inherited from DEG via
restriction. Combined with Step 3, the block restriction of DEG is
an equivalence by Lurie HA~4.8.2.11.
\end{proof}

\subsection{The Serre-Modular Sector}

\begin{lemma}[Khare--Wintenberger Serre's conjecture]
\label{v4-arith:w6-e:lem:KW-Serre}
\ClaimStatusProvedElsewhere. Every odd irreducible residual
representation
\(\bar\rho\colon G_{\mathbb{Q}}\to GL_{2}(\overline{\mathbb{F}}_{p})\)
is modular: there is a cuspidal eigenform \(f\) of weight
\(k(\bar\rho)\leq p+1\) and level
\(N(\bar\rho)\) coprime to \(p\) such that \(\bar\rho\cong\bar\rho_{f,p}\).
\end{lemma}

\begin{proof}[Source]
Khare--Wintenberger 2009 Ann.~Math.~169 (I), 170 (II) (full proof
of Serre's modularity conjecture). The argument proceeds via
compatible families and the Kisin 2009 modularity-lifting theorems.
\end{proof}

\begin{lemma}[Kisin modularity lifting under its hypotheses]
\label{v4-arith:w6-e:lem:Kisin-lifting}
\ClaimStatusProvedElsewhere. Let \(\bar\rho\) be odd irreducible and
residually modular with Serre weight and conductor
\((k(\bar\rho), N(\bar\rho))\). Let
\[
\rho\colon G_{\mathbb{Q}}\to GL_{2}(\overline{\mathbb{Q}}_{p})
\]
be a lift satisfying the local deformation, \(p\)-adic Hodge,
ramification, and Taylor--Wiles replacement hypotheses of Kisin 2009
J.~AMS~22 Thm~3.4. Then \(\rho\) is modular. Its local \(p\)-adic Hodge
type determines the corresponding Breuil--M{\'e}zard cycle.
\end{lemma}

\begin{proof}[Source]
Kisin 2009 J.~AMS~22 Thm~3.4 (modularity lifting from residual
modularity, without big-image hypothesis, using compatible families
plus Khare--Wintenberger). Gee 2011 Math.~Ann.~350 \S 2 supplies the
matching of prescribed \(p\)-adic Hodge types with Breuil--M{\'e}zard
cycles.
\end{proof}

\begin{theorem}[removal of Taylor--Wiles conditions on the Serre-modular sector]
\label{v4-arith:w6-e:thm:TW-removal-Serre}
\ClaimStatusNeedsVerification. For every prime \(p\geq 5\), every
\(\chi\in\mathrm{Hom}(\mathbb{Z}_{p}^{\times},k^{\times})\), and every
residual \(\bar\rho\in S^{\mathrm{Serre\text{-}res}}\), assume the
sectorial DEG generator comparison. Then the
restriction of \(\Phi^{\mathrm{DEG}}_{p,\chi}\) to the corresponding
block is an equivalence.
\end{theorem}

\begin{proof}
\textbf{Step 1: Serre modular generator.} By
\Cref{v4-arith:w6-e:lem:KW-Serre}, there is a cuspidal eigenform
\(f\) with \(\bar\rho\cong\bar\rho_{f,p}\). The block
\(\mathrm{Block}_{\bar\rho|_{G_{\mathbb{Q}_{p}}}}\) has the associated
compact projective \(\widetilde{\pi}_{f,p}\) (Colmez 2010 \S 7.2 for
the irreducible case; Pa{\v s}k{\=u}nas 2013 \S 2.6 for the
supersingular case).

\textbf{Step 2: modular lifts along the Breuil--M{\'e}zard
cycle.} By \Cref{v4-arith:w6-e:lem:Kisin-lifting}, every lift satisfying
Kisin's local and global hypotheses with the prescribed Hodge type
corresponds to a closed point of the Breuil--M{\'e}zard cycle
\(\mathcal{X}^{\mathrm{BM},k(\bar\rho)}_{p,\chi,\bar\rho}\). Gee 2011
identifies the cycle attached to this \(p\)-adic Hodge type. The
sectorial DEG generator comparison is the assertion that the projective
\(\widetilde{\pi}_{f,p}\) is sent to
\(\mathcal{O}_{\mathcal{X}^{\mathrm{BM},k(\bar\rho)}_{p,\chi,\bar\rho}}\).

\textbf{Step 3: compact-generator match.} The structure sheaves of
the Breuil--M{\'e}zard cycles indexed by Serre weights
\(\{\sigma\colon\sigma\text{ admissible for }\bar\rho\}\) form a
compact generator of
\(\mathrm{IndCoh}_{\mathrm{Nilp}}(\mathcal{X}^{\mathrm{EG},p,\chi,\bar\rho}_{GL_{2}})\)
(Emerton--Gee arXiv:1908.07185 \S 7, with the Serre weight
description via the Breuil--M{\'e}zard conjecture). The corresponding
Colmez projectives \(\{\widetilde{\pi}_{f,p}\}\) indexed by the same
weights form a compact generator of the block on the automorphic
side (Pa{\v s}k{\=u}nas 2013 \S 6).

\textbf{Step 4: equivalence by Lurie HA~4.8.2.11.} Fully-faithfulness
of DEG plus the assumed compact-generator match from Steps 1--3 gives
the equivalence on each Serre-residual block.
\end{proof}

\subsection{Assembly on the Union of Non-Taylor--Wiles Sectors}

\begin{theorem}[removal of Taylor--Wiles conditions on the union of sectors]
\label{v4-arith:w6-e:thm:TW-removal-union}
\ClaimStatusNeedsVerification. For every prime \(p\geq 5\), every central
character \(\chi\), and every residual
\(\bar\rho\in S^{\mathrm{dihedral}}\cup S^{\mathrm{nearly\text{-}ord}}\cup
S^{\mathrm{Serre\text{-}res}}\), assume the sectorial DEG generator
comparisons. Then the restriction of
\(\Phi^{\mathrm{DEG}}_{p,\chi}\) to the corresponding block is an
equivalence. Combined with
\Cref{v4-arith:w5-e:thm:f1a1-ordinary}, essential surjectivity of
\(\Phi^{\mathrm{DEG}}_{p,\chi}\) holds on the unramified ordinary
locus plus the union of the three non-Taylor--Wiles sectors only under
these comparison inputs.
\end{theorem}

\begin{proof}
Assemble \Cref{v4-arith:w6-e:thm:TW-removal-dihedral},
\Cref{v4-arith:w6-e:thm:TW-removal-nearly-ord}, and
\Cref{v4-arith:w6-e:thm:TW-removal-Serre}. Orthogonality of blocks
(each residual \(\bar\rho|_{G_{\mathbb{Q}_{p}}}\) indexes a separate
block of \(\mathrm{Mod}^{\mathrm{sm},\chi}\); the same residual-index
on the spectral side via
\Cref{v4-arith:deg-f1a:prop:spectral-blocks}) gives the direct-sum
assembly.
\end{proof}

\subsection{The Non-Solvable Non-Nearly-Ordinary Complement}

\begin{corollary}[the complementary sector]
\label{v4-arith:w6-e:cor:TW-removal-uncond}
\ClaimStatusNeedsVerification. Under the sectorial DEG generator
comparisons of \Cref{v4-arith:w6-e:thm:TW-removal-union}, the
complement to the covered non-Taylor--Wiles sectors is the sector
\(S^{\mathrm{non\text{-}solv}}\) of residuals
\(\bar\rho\) that are \emph{not} in
\(S^{\mathrm{dihedral}}\cup S^{\mathrm{nearly\text{-}ord}}\cup
S^{\mathrm{Serre\text{-}res}}\): genuinely non-solvable (projectively
\(A_{5}\)) and non-nearly-ordinary, with \(\bar\rho|_{G_{\mathbb{Q}_{p}}}\)
non-distinguished in the sense of (TW3) and not covered by
Skinner--Wiles.
\end{corollary}

\begin{proof}
Direct from \Cref{v4-arith:w6-e:def:non-TW-sectors} and
\Cref{v4-arith:w6-e:thm:TW-removal-union}. The non-solvable
residuals are exactly those with projective image
\(A_{5}\subset PGL_{2}(\overline{\mathbb{F}}_{p})\) or
non-trivial absolute Galois image beyond the solvable closure; by
Dickson's classification of finite subgroups of \(PGL_{2}\), the
projective images are cyclic, dihedral, \(A_{4}\), \(S_{4}\), \(A_{5}\),
or \(PSL_{2}(\mathbb{F}_{q})\) / \(PGL_{2}(\mathbb{F}_{q})\). The first
four are solvable and hence in \(S^{\mathrm{dihedral}}\); \(A_{5}\) is
exceptional; the large-image cases are in \(S^{\mathrm{Serre\text{-}res}}\)
by Khare--Wintenberger. The \(A_{5}\) sector plus its non-nearly-ordinary
extension is the residual non-solvable core.
\end{proof}

\begin{conjecture}[essential surjectivity on the non-solvable non-nearly-ordinary core]
\label{v4-arith:w6-e:conj:TW-removal-ns-residual}
\ClaimStatusConjectured. Essential surjectivity of
\(\Phi^{\mathrm{DEG}}_{p,\chi}\) on the non-solvable
non-nearly-ordinary core \(S^{\mathrm{non\text{-}solv}}\) holds if
Taylor's icosahedral potential automorphy, the relevant
Barnet-Lamb--Gee--Geraghty automorphy-lifting hypotheses, Emerton
patching, and the sectorial DEG generator comparison hold on this
block.
\end{conjecture}

\begin{remark}[the non-solvable residual]
\label{v4-arith:w6-e:rem:ns-residual-domain}
The \(A_{5}\) icosahedral case was handled by Taylor 2003 Pub.~IHES~98
Part~II in the classical Artin setting via potential automorphy and
solvable descent after a totally real base change. The general
BLGG 2012 J.~AMS~25 programme extends this to the case of Hilbert
modular forms over totally real fields, under its automorphy-lifting
hypotheses. Descent from a solvable extension to \(\mathbb{Q}\) and the
comparison of compact generators are additional hypotheses, supplied in
the form stated in \Cref{v4-arith:w6-e:conj:TW-removal-ns-residual}.
\end{remark}

\subsection{Independent Inputs for the Taylor--Wiles Sectors}

\paragraph{Sources.} Khare--Wintenberger 2009 proves
Serre's conjecture via compatible families. Arthur--Clozel 1989 is
solvable base change via the trace formula. Skinner--Wiles 1999 is
nearly-ordinary deformation theory. Taylor 2003 uses potential
automorphy with solvable-base-change descent. Barnet-Lamb--Gee--Geraghty 2012
is automorphy lifting via the Harris--Shepherd-Barron--Taylor
Calabi--Yau construction. These five sources have disjoint methods, and
all are independent of DEG arXiv:2603.26887 and
from the Kisin 2009 / Gee--Kisin 2014 / Emerton--Gee arXiv:1109.4226
sources used in \S\ref{v4-arith:w6-e:non-ord}.

\paragraph{Conventions.} Khare--Wintenberger uses the Serre weight
\(k(\bar\rho)\in\{2,\ldots,p+1\}\); Kisin uses dual-normalised
Hodge--Tate weights. These match via the standard shift
(Emerton--Gee arXiv:1908.07185 \S 3.1). Nearly-ordinary Hida-family
conventions (weight-space variable \(k\) versus Hecke operator
\(U_{p}\)) follow Skinner--Wiles 1999 \S 2.

\paragraph{Categorical framework.} The left-hand side is a block of
\(\mathrm{D}^{\mathrm{sph\text{-}fin}}([GL_{2}(\mathbb{Q}_{p})])_{\chi}\)
restricted to a specific residual \(\bar\rho|_{G_{\mathbb{Q}_{p}}}\);
the right-hand side is \(\mathrm{IndCoh}_{\mathrm{Nilp}}\) on the
corresponding component of the Emerton--Gee stack. Both are stable
presentable \(\infty\)-categories with compact generators indexed by
Serre weights or projective covers.

\paragraph{Local hypotheses.} Condition (TW1) (\(p\geq 5\)) is
retained. Conditions (TW2) (big image) and (TW3) (distinguishedness)
are replaced sector by sector: dihedral residuals are induced from a
quadratic extension; tetrahedral and octahedral residuals are governed
by Langlands--Tunnell; nearly-ordinary
residuals admit a nearly-ordinary lift; Serre-modular residuals are
covered by Khare--Wintenberger together with the hypotheses of Kisin's
lifting theorem. Dickson's classification separates the solvable,
icosahedral, and large-image projective cases.

\paragraph{Hecke compatibility.} The Hecke-eigenvalue match on each
sector is verified by the corresponding modularity theorem:
Arthur--Clozel \S 3 Thm 4 for the dihedral case;
\(R^{\mathrm{n\text{-}ord}}=T^{\mathrm{n\text{-}ord}}\) for the
nearly-ordinary case; Khare--Wintenberger plus Kisin modularity
lifting for the Serre-modular case.

\paragraph{Non-solvable comparison.} The three covered sectors use the
stated sectorial DEG generator comparisons. The non-solvable
non-nearly-ordinary sector is the comparison of
\Cref{v4-arith:w6-e:conj:TW-removal-ns-residual}.

\section{The Conditional Equivalence for \(\Phi^{\mathrm{DEG}}_{p,\chi}\)}
\label{v4-arith:w6-e:assembly}

\subsection{Union of the Strata}

\begin{theorem}[conditional essential surjectivity locus for \(\Phi^{\mathrm{DEG}}_{p,\chi}\)]
\label{v4-arith:w6-e:thm:f1a1-maximal-uncond}
\ClaimStatusNeedsVerification. For every prime \(p\geq 5\) and every central
character \(\chi\in\mathrm{Hom}(\mathbb{Z}_{p}^{\times},k^{\times})\),
assume the required DEG generator comparisons on the listed sectors.
Then
the DEG functor \(\Phi^{\mathrm{DEG}}_{p,\chi}\) is an equivalence on
the union of:
\begin{enumerate}
\item the unramified ordinary summand (\Cref{v4-arith:w5-e:thm:f1a1-ordinary},
under (TW1)--(TW4));
\item the potentially Barsotti--Tate non-ordinary summand
(\Cref{v4-arith:w6-e:thm:f1a1-pBT-nonord}, under (TW1)--(TW3));
\item the dihedral, nearly-ordinary, and Serre-modular non-Taylor--Wiles
sectors (\Cref{v4-arith:w6-e:thm:TW-removal-union}, under (TW1) only).
\end{enumerate}
The deformation and automorphy sources identify the relevant strata and
generators. The categorical assertion is the compatibility of these
generators with \(\Phi^{\mathrm{DEG}}_{p,\chi}\).
\end{theorem}

\begin{proof}
The three stated theorems apply on pairwise orthogonal
residual-index blocks. Their direct sum is therefore an equivalence on
the union.
\end{proof}

\subsection{Boundary of the Non-Ordinary and Non-Taylor--Wiles Comparison}

\begin{proposition}[residual decomposition for essential surjectivity on the non-ordinary and non-Taylor--Wiles loci]
\label{v4-arith:w6-e:prop:F1-table-refined}
\ClaimStatusProvedHere. After Chapter~\ref{v4-arith:w5-e:chapter}
and this chapter, the extension of
\(\Phi^{\mathrm{DEG}}_{p,\chi}\) from the ordinary, pBT, solvable,
nearly-ordinary, and Serre-modular strata to the full non-ordinary and
non-Taylor--Wiles locus is equivalent to the following two generator
comparisons:
\begin{enumerate}
\item Supersingular core
(\Cref{v4-arith:w6-e:conj:f1a1-ss-residual}): essential surjectivity
on the genuinely supersingular sub-locus at \(p\geq 5\).
\item Non-solvable non-nearly-ordinary core
(\Cref{v4-arith:w6-e:conj:TW-removal-ns-residual}): essential
surjectivity on the projectively \(A_{5}\) and non-nearly-ordinary
residual sector.
\end{enumerate}
\end{proposition}

\begin{proof}
The non-ordinary locus is
\(\mathcal{X}^{\mathrm{pBT}}_{\bar\rho,\chi}\cup
\mathcal{X}^{\mathrm{ss}}_{\bar\rho,\chi}\). The pBT part is governed by
\Cref{v4-arith:w6-e:thm:f1a1-pBT-nonord}; its complement is
\Cref{v4-arith:w6-e:conj:f1a1-ss-residual}. The non-Taylor--Wiles
residuals decompose as
\[
S^{\mathrm{dihedral}}\cup S^{\mathrm{nearly\text{-}ord}}\cup
S^{\mathrm{Serre\text{-}res}}\cup S^{\mathrm{non\text{-}solv}}.
\]
The first three sectors are governed by
\Cref{v4-arith:w6-e:thm:TW-removal-union}; the last sector is
\Cref{v4-arith:w6-e:conj:TW-removal-ns-residual}.
\end{proof}

\subsection{Interaction with the Full Essential Surjectivity Problem}

\begin{proposition}[essential surjectivity after the ordinary and non-ordinary strata]
\label{v4-arith:w6-e:prop:F1-table-combined}
\ClaimStatusProvedHere. Combining
\Cref{v4-arith:w5-e:prop:F1-table-refined} with
\Cref{v4-arith:w6-e:prop:F1-table-refined}, essential surjectivity of
\(\Phi^{\mathrm{DEG}}_{p,\chi}\) on the full locus is equivalent to the
following five comparisons:
\begin{enumerate}
\item Supersingular core.
\item Non-solvable non-nearly-ordinary core.
\item \textbf{F1.a.2.class-match}
(\Cref{v4-arith:f1a3-BH:conj:Paskunas-wild-ext}).
\item \textbf{\(\mathrm{L}_{\mathrm{ACD}}\).Eisenstein}
(\Cref{v4-arith:w5-e:conj:LACD-Eisenstein}).
\item \textbf{BC-diamond-glue}
(\Cref{v4-arith:bcglue-core}).
\end{enumerate}
\end{proposition}

\begin{proof}
Direct combination of \Cref{v4-arith:w5-e:prop:F1-table-refined}
and \Cref{v4-arith:w6-e:prop:F1-table-refined}.
\end{proof}

\section{Consequences}
\label{v4-arith:w6-e:summary}

The non-ordinary locus is the union of the pBT stratum and the
supersingular complement. The pBT stratum is controlled by Kisin,
Gee--Kisin, Emerton--Gee, Colmez, and the pBT DEG generator comparison.
The supersingular complement is the separate generator comparison of
\Cref{v4-arith:w6-e:conj:f1a1-ss-residual}.

The failure of (TW2) or (TW3) decomposes into solvable projective-image,
nearly-ordinary, Serre-modular, and non-solvable non-nearly-ordinary
sectors. The first three are controlled by
\Cref{v4-arith:w6-e:thm:TW-removal-union}; the last is the comparison
\Cref{v4-arith:w6-e:conj:TW-removal-ns-residual}. Thus the categorical
content is concentrated in explicit compact-generator comparisons, not
in the deformation or automorphy theorems alone.
